\section{Question 3: Etude de validité}

\subsection{Droit applicable}

La loi applicable est l'article L614-12 CPI renvoyant
à l'article 138 CBE.\\

\textbf{Conclusion: le motif de l'article 138(1)a) CBE
est analysé.}

\subsection{Divulgations opposables}

Le brevet EP n'a pas de priorité. Ainsi, les objets
définis par R1 et R2 bénéficient de la date de dépôt.
Ainsi, la date effective est le 22/06/2013
pour toutes les revendications.

\begin{itemize}
    \item D1 est une publication de 1996;
    \item D2 est une demande de brevet FR publiée en 2005;
    \item D3 est une demande de brevet EP publiée le 12/09/2011.
\end{itemize}
Donc, D1, D2 et D3 sont publiées avant la date effective.\\

\textbf{Conclusion: D1, D2 et D3 sont des états de la technique \epcart{54}(2) CBE de la demande de brevet EP.}\\

\subsection{Nouveauté}

\subsubsection{Nouveauté par rapport à D1}

D1 divulgue un dispositif 9 pour l’écartement automatique
de deux banches lors de leur levage (figure 1), comprenant :

\begin{itemize}
\item un compas 10 formé par une tête 11 et deux branches 12, 13.\\
\end{itemize}

D1 ne divulgue pas un dispositif comprenant:
\begin{itemize}
    \item les branches
12, 13 comprenant une pièce de liaison 15 apte à être montée de façon réversible sur la
face supérieure d’une banche,
    \item un montant vertical 18 comprenant une extrémité supérieure 19 pour la préhension du
dispositif et une extrémité inférieure reliée aux branches 12, 13 du compas 10 par
l’intermédiaire de bielles 21, 22 respectives, le montant vertical 18 étant arrangé
coulissant au travers de la tête 11 du compas 10.\\
\end{itemize}

R2 dépend de R1. D2 ne divulgue pas R2.\\

\textbf{Conclusion: R1 et R2 sont nouvelles par rapport à D1.}

% \subsubsection{Nouveauté par rapport à D2}

% \subsubsection{Nouveauté par rapport à D3}

% \subsection{Activité inventive}